\documentclass[12pt]{article}
\usepackage[utf8]{inputenc}
\usepackage[T1]{fontenc}
\usepackage{hyperref}
\title{Makwe Language Resources for K-12 Curriculum Development}
\author{Compiled by Claude Sonnet 5.5 (claude-sonnet-5-5) for Andr\'as Kornai}
\date{}
\begin{document}
\maketitle

\section{Introduction}
Makwe (Kimakwe, Macue, Maraba in Tanzania; ISO 639-3 ymk) is a Bantu language of the Cabo Delgado coast of Mozambique and Mtwara Region of Tanzania. It is lexically close to Swahili (about 60\%) but not intelligible with it. Estimates of speaker numbers vary widely between sources. This report rests on a single web search, and little material specific to schooling was found.

\section{Primary Educational Resources}
No primers, textbooks or school programmes in Makwe were found in the search results.

\section{Linguistic and Reference Materials}
\begin{itemize}
\item Makwe language overview (Wikipedia), which mentions the 2008 grammar by Maud Devos: \url{https://en.wikipedia.org/wiki/Makwe_language}
\item Devos, Maud 2004, bibliographic record in Glottolog: \url{https://glottolog.org/resource/reference/id/87970}
\item A PDF at dwc.knaw.nl appeared in the results. I did not check its content, so its relevance is unconfirmed: \url{https://dwc.knaw.nl/DL/publications/PU00010055.pdf}
\end{itemize}

\section{Religious and Community Resources}
\begin{itemize}
\item Makwe people (Wikipedia): \url{https://en.wikipedia.org/wiki/Makwe_people}
\item Joshua Project language page: \url{https://www.joshuaproject.net/index.php/languages/ymk}
\item Joshua Project, Makwe in Tanzania: \url{https://joshuaproject.net/index.php/people_groups/13429/TZ}
\item Afrigo article on the Makwe of Tanzania and Mozambique: \url{https://afrigo.org/articles/people-group-the-makwe-of-tanzania-and-mozambique/}
\end{itemize}

\section{Digital and Audio Resources}
No audio, video or corpus resources were found.

\section{Curriculum Development Context}
The search found no evidence of Makwe-medium instruction. A curriculum would have to start from the Devos grammar and from new community-based documentation. Swahili and Portuguese are the likely school languages in the region, but this was not confirmed by the search.
\end{document}